\documentclass{article}

\usepackage{amsmath,amssymb}
\usepackage{xurl}
\usepackage[hidelinks]{hyperref}
\setlength{\emergencystretch}{2em}

\input{command}

\title{Hypothesis Boundaries for the Indecomposability\\
of the Breuer--Hall Map (Wolf Example~3.1)}
\author{}
\date{2026-08-23}

\begin{document}
\maketitle
\gapnote{local-correction}{resolved}

\section{Source statement}

Let $U\in\MD$ be antisymmetric, so $U^{\mathsf T}=-U$, and suppose that it is
a contraction:
\begin{align}
    U^\dagger U
        &\leq\Id.
    \label{eq:wolf_breuer_hall_even_dim_restriction_contraction}
\end{align}
Wolf's Chapter~3, Example~3.1 defines the Breuer--Hall map by
\begin{align}
    T_{\mathrm{BH}}(X)
        &=\tr(X)\Id-X-UX^{\mathsf T}U^\dagger.
    \label{eq:wolf_breuer_hall_even_dim_restriction_breuer_hall_map}
\end{align}
The source identifies this formula as Equation~(3.19)
\cite[Chapter~3, Example~3.1]{Wolf2012Quantum}.\footnote{Local source:
\path{Notes/WolfNoteTexSource/ch03_positive_not_completely.tex}, lines
344--355. Formalization issue: \ghissue{16}.}

Wolf states that antisymmetric unitaries exist when $d$ is even and satisfy
``$U^\dagger=U^{-1}=-\overline U$''. The source then states, without an
explicit dimension hypothesis, that $T_{\mathrm{BH}}$ is not decomposable
when $U$ is an antisymmetric unitary \cite[Example~3.1]{Wolf2012Quantum}.
Its positivity and indecomposability assertions therefore use different
hypotheses: positivity allows the contraction condition
(\ref{eq:wolf_breuer_hall_even_dim_restriction_contraction}), whereas
indecomposability requires
\begin{align}
    U^\dagger U
        &=\Id.
    \label{eq:wolf_breuer_hall_even_dim_restriction_unitarity}
\end{align}

The formalization preserves this distinction. It proves positivity for
antisymmetric contractions and indecomposability for antisymmetric unitaries
when $d>2$. Indecomposability does not extend to all antisymmetric
contractions.

\section{The dimension-two boundary}

Every antisymmetric $2\times2$ matrix is a scalar multiple of
\begin{align}
    J
        &=\begin{pmatrix}
            0 & 1\\
            -1 & 0
          \end{pmatrix}.
    \label{eq:wolf_breuer_hall_even_dim_restriction_symplectic_matrix}
\end{align}
Thus every antisymmetric unitary in dimension two has the form
$e^{i\theta}J$, and the phase does not affect the twisted-transpose term.
For
\begin{align}
    X
        &=\begin{pmatrix}
            a & b\\
            c & d
          \end{pmatrix},
    \label{eq:wolf_breuer_hall_even_dim_restriction_generic_two_matrix}
\end{align}
direct multiplication gives
\begin{align}
    JX^{\mathsf T}J^\dagger
        &=\begin{pmatrix}
            d & -b\\
            -c & a
          \end{pmatrix},
    \label{eq:wolf_breuer_hall_even_dim_restriction_twisted_transpose_matrix}\\
    JX^{\mathsf T}J^\dagger
        &=\tr(X)\Id-X.
    \label{eq:wolf_breuer_hall_even_dim_restriction_twisted_transpose_reduction}
\end{align}
Substitution into
(\ref{eq:wolf_breuer_hall_even_dim_restriction_breuer_hall_map}) yields
\begin{align}
    T_{\mathrm{BH}}(X)
        &=0
    \label{eq:wolf_breuer_hall_even_dim_restriction_dimension_two_zero}
\end{align}
for every $X\in M_2(\C)$. The zero map is completely positive and therefore
decomposable. Hence $d>2$ is necessary.

An antisymmetric unitary can exist only in even dimension. Indeed,
$U^{\mathsf T}=-U$ implies
\begin{align}
    \det U
        &=\det(U^{\mathsf T}),
    \label{eq:wolf_breuer_hall_even_dim_restriction_determinant_transpose}\\
    \det U
        &=\det(-U),
    \label{eq:wolf_breuer_hall_even_dim_restriction_determinant_antisymmetry}\\
    \det U
        &=(-1)^d\det U.
    \label{eq:wolf_breuer_hall_even_dim_restriction_determinant_parity}
\end{align}
Since unitarity gives $\det U\neq0$, the dimension is even. The corrected
unitary range is therefore $d\geq4$.

\section{Why contractivity does not imply indecomposability}

The zero matrix is an antisymmetric contraction in every dimension:
\begin{align}
    0^{\mathsf T}
        &=-0,
    \label{eq:wolf_breuer_hall_even_dim_restriction_zero_antisymmetric}\\
    0^\dagger0
        &\leq\Id.
    \label{eq:wolf_breuer_hall_even_dim_restriction_zero_contraction}
\end{align}
At $U=0$, the Breuer--Hall map becomes the first reduction map:
\begin{align}
    T_{\mathrm{BH},0}(X)
        &=\tr(X)\Id-X,
    \label{eq:wolf_breuer_hall_even_dim_restriction_zero_parameter_map}\\
    T_{\mathrm{BH},0}(X)
        &=T_1(X).
    \label{eq:wolf_breuer_hall_even_dim_restriction_reduction_map}
\end{align}

Let $\theta$ denote transposition and $F$ the swap operator. With the
normalization used in the formal development, the Choi matrix of
$T_1\circ\theta$ is
\begin{align}
    \tau(T_1\circ\theta)
        &=d^{-1}(\Id-F),
    \label{eq:wolf_breuer_hall_even_dim_restriction_reduction_choi}\\
    \tau(T_1\circ\theta)
        &\geq0.
    \label{eq:wolf_breuer_hall_even_dim_restriction_reduction_choi_positive}
\end{align}
The second line holds because $\Id-F$ is twice the orthogonal projector onto
the antisymmetric subspace. The Choi criterion shows that $T_1\circ\theta$ is
completely positive. Therefore $T_1$ is completely copositive and
decomposable.

Consequently, even when $d>2$, the Breuer--Hall map need not be
indecomposable over all antisymmetric contractions. This counterexample does
not classify the nonunitary contractions for which the map is
indecomposable, and the formalization makes no such claim.\footnote{Formal
declarations:
\leanid{Matrix.breuerHallMap\_zero} in
\doclink{QICLean/Channel/BreuerHallMap};
\leanid{ChoiJamiolkowski.choiMatrix\_reductionMap\_one\_comp\_transpose},
\leanid{Matrix.reductionMap\_one\_isCompletelyCopositiveMap}, and
\leanid{Matrix.reductionMap\_one\_isDecomposablePositiveMap} in
\doclink{QICLean/Channel/PositiveExamples}; and
\leanid{Matrix.breuerHallMap\_contraction\_not\_universally\_indecomposable}
in \doclink{QICLean/Channel/BreuerHallIndecomposable}.}

\section{The corrected statement}

If $d>2$ and $U$ is an antisymmetric unitary on $\C^d$, then
$T_{\mathrm{BH}}$ is an indecomposable positive map.\footnote{Formal
declarations:
\leanid{Matrix.breuerHallMap\_isIndecomposablePositiveMap} and
\leanid{Matrix.breuerHallMap\_not\_isDecomposablePositiveMap} in
\doclink{QICLean/Channel/BreuerHallIndecomposable}; positivity is
\leanid{Matrix.breuerHallMap\_isPositiveMap} in
\doclink{QICLean/Channel/BreuerHallMap}.}

The source gives no proof. The formalized argument uses PPT detection. Let
$F$ be the swap operator on $\C^d\otimes\C^d$, and let $\ket{\Psi}$ be the
$U$-twisted maximally entangled vector with coordinates
$\Psi(i,k)=U(i,k)$. Define the unnormalized projector and witness state by
\begin{align}
    P_0
        &=\ket{\Psi}\bra{\Psi},
    \label{eq:wolf_breuer_hall_even_dim_restriction_twisted_projector}\\
    \rho
        &=(\Id+F)+P_0.
    \label{eq:wolf_breuer_hall_even_dim_restriction_witness_state}
\end{align}
Both terms in
(\ref{eq:wolf_breuer_hall_even_dim_restriction_witness_state}) are positive
semidefinite. Let $P_1$ denote the untwisted maximally entangled projector.
Antisymmetry and unitarity give
\begin{align}
    (\Id+F)^{T_1}
        &=\Id+P_1,
    \label{eq:wolf_breuer_hall_even_dim_restriction_swap_partial_transpose}\\
    P_0^{T_1}
        &=(U^\dagger\otimes\Id)(\Id+F)(U\otimes\Id)-\Id.
    \label{eq:wolf_breuer_hall_even_dim_restriction_twisted_partial_transpose}
\end{align}
Therefore
\begin{align}
    \rho^{T_1}
        &=P_1+(U^\dagger\otimes\Id)(\Id+F)(U\otimes\Id),
    \label{eq:wolf_breuer_hall_even_dim_restriction_witness_partial_transpose}\\
    \rho^{T_1}
        &\geq0.
    \label{eq:wolf_breuer_hall_even_dim_restriction_witness_ppt}
\end{align}
Thus $\rho$ is a PPT state.

Expanding the tensor indices and using $U^\dagger U=\Id$ gives, for every
bipartite matrix $\sigma$,
\begin{align}
    \tr\!\left[(T_{\mathrm{BH}}\otimes\operatorname{id})(\sigma)P_0\right]
        &=\tr(\sigma)-\tr(\sigma P_0)-\tr(F\sigma).
    \label{eq:wolf_breuer_hall_even_dim_restriction_pairing_identity}
\end{align}
For $\sigma=\rho$, antisymmetry gives
$F\ket{\Psi}=-\ket{\Psi}$, and the required traces are
\begin{align}
    \tr(\rho)
        &=d^2+2d,
    \label{eq:wolf_breuer_hall_even_dim_restriction_witness_trace}\\
    \tr(\rho P_0)
        &=d^2,
    \label{eq:wolf_breuer_hall_even_dim_restriction_projector_pairing}\\
    \tr(F\rho)
        &=d^2.
    \label{eq:wolf_breuer_hall_even_dim_restriction_swap_pairing}
\end{align}
Substitution into
(\ref{eq:wolf_breuer_hall_even_dim_restriction_pairing_identity}) yields
\begin{align}
    \tr\!\left[(T_{\mathrm{BH}}\otimes\operatorname{id})(\rho)P_0\right]
        &=d(2-d),
    \label{eq:wolf_breuer_hall_even_dim_restriction_negative_detection_value}\\
    \tr\!\left[(T_{\mathrm{BH}}\otimes\operatorname{id})(\rho)P_0\right]
        &<0
        \qquad(d>2).
    \label{eq:wolf_breuer_hall_even_dim_restriction_negative_detection}
\end{align}
Hence $(T_{\mathrm{BH}}\otimes\operatorname{id})(\rho)$ is not positive
semidefinite. Decomposable maps preserve positivity on PPT
states.\footnote{Formal declarations:
\leanid{IsDecomposablePositiveMap.tensorMapId\_posSemidef\_of\_isPPT} and
\leanid{not\_isDecomposablePositiveMap\_of\_isPPT\_not\_tensorMapId\_posSemidef}
in \doclink{QICLean/Channel/DecomposablePPT}.}
Therefore $T_{\mathrm{BH}}$ is not decomposable.

\section{Conclusion}

Wolf's positivity statement applies to antisymmetric contractions, while the
indecomposability statement applies only to antisymmetric unitaries
\cite[Example~3.1]{Wolf2012Quantum}. In the unitary case, the dimension
correction is necessary: the formal theorem assumes $d>2$, equivalently even
$d\geq4$ when an antisymmetric unitary exists, and the map vanishes at $d=2$.
The unitary hypothesis is also a genuine boundary for the stated conclusion,
because the antisymmetric contraction $U=0$ gives the completely copositive,
decomposable reduction map $T_1$. The formalization records the corrected
unitary theorem and this counterexample without asserting a classification for
general contractions.

\bibliographystyle{alpha}
\bibliography{references}

\end{document}